\subsection{معرفی و انگیزه}
این ماژول در چارچوب پژوهش حاضر پیاده‌سازی و ارزیابی شده و فاقد مقاله مستقل جداگانه در ادبیات تحقیق است.

در مدل‌های زبانی حجیم (\lr{LLMs})، ابعاد فعال‌سازی و تعداد نورون‌ها در لایه‌های پیش‌خور و توجه چندسر عامل اصلی اشغال حافظه گرافیکی و محاسبات تکراری در مرحله استنتاج هستند. روش‌های هرس وزن سنتی غالباً فرض می‌کنند که نورون‌های با وزن‌های کوچک‌تر، تأثیر کمتری بر خروجی دارند، در حالی که در عمل خروجی لایه‌ها تحت تأثیر ویژگی‌های ورودی با مقیاس‌های متغیر قرار دارد و انرژی خروجی واقعی نورون معیار بسیار دقیق‌تری برای تعیین اهمیت آن است.

ماژول فشرده‌سازی \lr{FLAT-LLM (Eigen-Decomposition / Energy-Aware Structured Pruning)} بر مبنای تحلیل کوواریانس و انرژی فعال‌سازی‌های خروجی نورون‌ها طراحی شده است. ایده بنیادین در این ماژول، تخمین مستقیم ماتریس کوواریانس خروجی و قطری‌سازی انرژی کانال‌های خروجی با استفاده از نمونه‌های کالیبراسیون است. با اولویت‌بندی کانال‌ها بر اساس مقدار انرژی خروجی تجربی، ماژول لایه‌ها را بدون از بین بردن قوام اطلاعاتی مدل سبک‌سازی کرده و امکان هرس ساختاریافته نورونی یا اعمال تنکی سخت‌افزاری $2:4$ را فراهم می‌سازد.
